%% sections/06-conclusion.tex

\section{Conclusion}
\label{sec:conclusion}

We have described a pipeline for converting session-level behavioral surprises into
ratified player-style design patterns, applied across 49 AI-simulated TRPG sessions
in 7 complete campaigns. The pipeline promoted 7 player styles, validated them across
12 research hypotheses spanning campaigns with professional motivations, duty-based
institutional structures, and minimal personal backstory, and confirmed that all 7
styles are design-enabled rather than archetype-specific.

The central practical contribution is a set of design prescriptions. A designer who
wants sheet-deep-reader behavior writes denser, more specific character sheets with
named losses, local-origin notes, and institutional-training histories. A designer
who wants craft-witness behavior writes PC sheets with named daily rituals that specify
a physical gesture or attention habit. A designer who wants act-without-announcement
instances structures campaigns to provide a naming event in an earlier session and a
scene in a later session where the named thing is available to act from. These
prescriptions are learnable, portable, and independent of the motivational register
of the campaign.

The trigger-type shift finding (RV11) extends the prescriptions to temporal design:
in long campaigns, the relevant trigger material for later-session instances is not
scene design but relationship history accumulated in prior sessions. The implication
is that designers working on multi-session campaigns should plan not only for
instance-enabling scenes but for the naming events that will load the later-session
instances.

The partial confirmation of RV8 --- the one incomplete result --- is also informative.
Verbal-documentation-arc may require oral-practice PC design to fire reliably, and
the formal-reporting variant may be a weaker trigger. This points toward a refinement
of the style definition: the style may be more precisely a \textit{spoken-documentation-arc}
(oral or performative) rather than a general documentation-practice arc. Future work
should test this refinement across campaigns designed with formal-reporting PC sheets.

The most significant open question is human generalization. Every result in this
paper was produced by an AI system playing PCs according to explicit Playstyle
Heuristics. Whether the same styles fire when human players play the same PC sheets
--- and whether the instance counts are higher or lower --- is an empirical question
that the current corpus cannot answer. The most probable outcome is that human play
produces the same styles at higher instance rates (human players improvise beyond
heuristics) and introduces additional noise (player personality bleeding through).
What the AI corpus provides is a design-theory foundation: the styles are demonstrably
enabled by character-sheet design, independently of who is playing. That foundation
is what makes them useful as design tools.

\bigskip
\noindent\textbf{Acknowledgments.} The Marathon D\&D Workshop corpus was developed using
the Marathon workshop framework and the Claude AI system (Anthropic, 2025--2026).
